Each criterion is recomputed for every retained conditional-quantile draw. Table entries report the posterior mean and equal-tailed 95\% credible interval of the resulting draw-wise criterion. Fit RMSE measures recovery of the true conditional quantile over the 500-observation training sample; forecast MAE and check loss average the 1,000 rolling-origin and horizon evaluations used in the point comparison. The intervals condition on the simulated data, evaluation design, and criterion-specific fitted model. Variational Bayes intervals are approximate.

The exDQLM entries use exdqlm 1.1.1; the DQLM and both Q--DESN entries use their stated specifications. In rolling-origin exDQLM MCMC evaluation, the initial parameter posterior is retained while the filtering distribution is updated sequentially with responses available at each origin. The observation-error mean and variance used in this update are averaged over paired retained draws of the scale and asymmetry parameters before the updated state is propagated over the requested forecast horizons. Point estimates and posterior intervals use distinct estimators: point estimates summarize a single fixed path for VB and chain-level fixed paths for MCMC, whereas interval centers are posterior means of draw-wise criteria. The two summaries are therefore interpreted separately.

5 of the 108 displayed MCMC summaries are marked by daggers to indicate diagnostic cautions; the supplement provides details.
